\section{Finite diffuse data and an intrinsic transposition Dirichlet problem}
\label{sec:r11-diffuse-dirichlet}

Retain the arbitrary original background and the analytic domains
$(B,\Omega,V)$ of Section~\ref{sec:r11-affinoid-energy}.  Thus
$S$ has degree $(n-1,n-1)$, $F=V\cap\operatorname{supp}S$,
$\lambda=D(hS)$ is faithful, and
$H=\mathscr D_0(V)$ is a regular strongly local coercive part form.
The reference measure $\lambda|_F$ is finite because $V\Subset\Omega$.
Write $c_V$ for the transient $0$-capacity of this part form.
Its usual outer definition agrees, on quasi-open sets, with
\[
 c_V(A)=\inf\{\mathcal E(v,v):v\in H,\ 
                              \widetilde v\geq1\text{ quasi-everywhere on }A\}.
 \tag{MD1}
\]
The tilde denotes the energy-quasicontinuous representative.
Since $\mathcal E\leq\mathcal E_1\leq3\mathcal E$,
\[
                    c_V\leq\operatorname{cap}_{H,1}\leq3c_V.
 \tag{MD2}
\]
The two capacities therefore have identical polar sets.
They are capacities on the original support space $F$, with
its ordinary subspace topology; they are not identified with
relative BT capacity.

\begin{definition}[Finite energy-diffuse data]
\label{def:r11-energy-diffuse-data}
Let $\mathcal M_{\mathrm{diff}}(F)$ be the finite signed Radon
measures $\mu$ on $F$ for which $|\mu|$ vanishes on every
$c_V$-polar Borel set.  This is a condition on the specified
background and domain, not on an unknown solution.
\end{definition}

\noindent\textbf{Status: definition.}

\begin{lemma}[The finite-energy nest input]
\label{lem:r11-finite-energy-nest}
For every $\mu\in\mathcal M_{\mathrm{diff}}(F)$ there are increasing
Borel sets $F_j$ such that
$|\mu|(F\setminus\bigcup_jF_j)=0$ and
$\tau_j={\bf1}_{F_j}|\mu|$ has finite energy integral:
\[
             \int_F|\widetilde v|\,d\tau_j
                         \leq C_j\,\mathcal E(v,v)^{1/2}
                                  \qquad(v\in H).
 \tag{MD3}
\]
In particular $\mu_j={\bf1}_{F_j}\mu$ converges to $\mu$ in
total variation and has a finite-energy variation.
\end{lemma}

\noindent\textbf{Status: \texttt{EXTERNAL\_KNOWN}}.
The precise input is Klimsiak--Rozkosz,
\emph{On the structure of bounded smooth measures associated with a
quasi-regular Dirichlet form}, arXiv:1410.4927v2,
Lemma~2.1 (pp.~2--3).  It asserts a finite $0$-energy nest
for every smooth measure of a transient quasi-regular form.
The original form here is regular, hence quasi-regular, and
coercivity proves transience.  A finite Radon measure charging
no polar sets is smooth: any compact form nest has finite
measure on its members and has polar complement.  Apply the
lemma to $|\mu|$.  The claimed total-variation convergence
then follows by monotone convergence.
No existence theorem for the original current equation is
included in this external input.

\subsection{Quantitative estimates before passing to measure data}

For a finite signed Radon measure $\sigma$ whose variation has
finite energy integral, let $G\sigma\in H$ be its Riesz potential:
\[
               \mathcal E(G\sigma,v)=\int_F\widetilde v\,d\sigma
                                 \qquad(v\in H).
 \tag{MD4}
\]
This is the R6 finite-energy result applied on the analytic
domains now covered by Section~\ref{sec:r11-affinoid-energy}.
For $k>0$ set $T_k(s)=\max(-k,\min(k,s))$.

\begin{lemma}[Mass controls integrability and all truncations]
\label{lem:r11-finite-energy-mass-estimates}
Under the preceding assumptions on $\sigma$, put $w=G\sigma$.
Then
\[
 \|w\|_{L^1(\lambda)}\leq2\|\sigma\|_{\mathrm{TV}},\qquad
 \mathcal E(T_k w,T_k w)\leq k\|\sigma\|_{\mathrm{TV}},\qquad
 c_V(\{|\widetilde w|>k\})\leq\frac{\|\sigma\|_{\mathrm{TV}}}{k}.
 \tag{MD5}
\]
The constants are independent of the finite-energy norm of $\sigma$.
\end{lemma}

\noindent\textbf{Status: \texttt{NEW\_PROVED}} as the uniform
original-background mass estimate needed for the larger solution class.
Only the frozen positivity and coercive barrier are used.

\begin{proof}
Positive finite-energy measures integrate canonical energy
representatives.  To recall the justification, their core norm
bound and the compact-core capacity formula show that they
vanish on polar compact sets, hence on polar Borel sets by Radon
regularity.  For a core approximation $v_j\to v$ in $H$,
the same norm bound applied to smooth approximations of
$|v_j-v_l|$ makes it Cauchy in $L^1$ of the measure.
A subsequence converges quasi-everywhere, so its $L^1$ limit is
the canonical representative $\widetilde v$.
This proves the integration interpretation in (MD4).

Order preservation of $G$ gives
$|G\sigma|\leq G|\sigma|$ quasi-everywhere.  With the torsion
vector $z=G1$, whose representative satisfies $0\leq z\leq2$,
symmetry yields
\[
 \int_FG|\sigma|\,d\lambda
       =\mathcal E(G|\sigma|,z)
       =\int_F\widetilde z\,d|\sigma|
       \leq2\|\sigma\|_{\mathrm{TV}}.
 \tag{MD6}
\]
This proves the first estimate.

The normal contraction $P(s)=s-2T_k(s)$ has Lipschitz constant
one and $P(0)=0$.  The Markov contraction inequality therefore
gives
$\mathcal E(w-2T_kw,w-2T_kw)\leq\mathcal E(w,w)$.
Expanding proves
$\mathcal E(T_kw,T_kw)\leq\mathcal E(w,T_kw)$.
Equation (MD4), tested with $T_kw$, bounds its right side by
$k\|\sigma\|_{\mathrm{TV}}$.  The vector $|T_kw|/k$ is
admissible for the outer-capacity estimate of the quasi-open
set $\{|\widetilde w|>k\}$, giving the last inequality.
\end{proof}

\begin{definition}[The transposition boundary condition]
\label{def:r11-transposition-boundary}
A function $u\in L^1(F,\lambda)$ satisfies the zero transposition
Dirichlet problem for $\mu\in\mathcal M_{\mathrm{diff}}(F)$ if
\[
             \int_Fu f\,d\lambda
                   =-\int_F\widetilde{Gf}\,d\mu
                        \qquad(f\in L^\infty(F,\lambda)).
 \tag{MD7}
\]
Here $Gf=G(f\lambda)\in H$, and its canonical representative is
bounded by $2\|f\|_\infty$.  The right side is therefore defined
for every datum under consideration.
\end{definition}

\noindent\textbf{Status: definition.}
This is an intrinsic energy boundary condition through the
part resolvent.  It does not assert an ordinary pointwise
boundary limit.

\begin{theorem}[All-original-current Dirichlet solutions for finite diffuse data]
\label{thm:r11-diffuse-transposition-dirichlet}
Every $\mu\in\mathcal M_{\mathrm{diff}}(F)$ has a unique
$L^1(F,\lambda)$ solution $u_\mu$ of (MD7).  The solution has
a canonical finite energy-quasicontinuous representative
outside a polar set, and
\[
       T_k(u_\mu)\in H,\qquad
       \mathcal E(T_ku_\mu,T_ku_\mu)\leq k\|\mu\|_{\mathrm{TV}}.
 \tag{MD8}
\]
Its coefficient-defined original current satisfies
\[
                         D(u_\mu S)=\mu\quad\hbox{on }V
 \tag{MD9}
\]
against every original compact smooth scalar test, including
all actual moment refinements.
The solution map is linear, is order reversing on positive data,
and for $\mu,\nu\in\mathcal M_{\mathrm{diff}}(F)$ satisfies
\[
 \begin{split}
 \|u_\mu-u_\nu\|_{L^1(\lambda)}
                   &\leq2\|\mu-\nu\|_{\mathrm{TV}},\\
 \mathcal E\bigl(T_k(u_\mu-u_\nu),T_k(u_\mu-u_\nu)\bigr)
                   &\leq k\|\mu-\nu\|_{\mathrm{TV}},\\
 c_V(\{|\widetilde u_\mu-\widetilde u_\nu|>k\})
                   &\leq\|\mu-\nu\|_{\mathrm{TV}}/k.
 \end{split}
 \tag{MD10}
\]
For positive $\mu$, one has $u_\mu\leq0$.
When $\mu$ already has finite energy, this is exactly the
frozen R6 energy solution.
\end{theorem}

\noindent\textbf{Status: \texttt{NEW\_PROVED}} as the
all-original-current extension to finite diffuse data without a
finite-energy norm assumption.  The abstract exhaustion lemma
above is the stated \texttt{EXTERNAL\_KNOWN} input.
This theorem does not assert ambient classical psh membership
or a nonlinear diagonal Monge--Ampere equation.

\begin{proof}
Choose $\mu_j$ from Lemma~\ref{lem:r11-finite-energy-nest} and
put $u_j=-G\mu_j$.  Differences have finite-energy variation,
so (MD5) gives
\[
 \|u_j-u_l\|_1\leq2\|\mu_j-\mu_l\|_{\mathrm{TV}},\qquad
 c_V(\{|\widetilde u_j-\widetilde u_l|>k\})
       \leq\|\mu_j-\mu_l\|_{\mathrm{TV}}/k.
 \tag{MD11}
\]
Hence $u_j$ converges in $L^1(\lambda)$ to some $u$ and is
Cauchy in energy capacity.  A subsequence with summable
successive capacity errors has a finite limit outside a polar
set: choose summable level thresholds and cover the corresponding
successive gap sets by open sets with summable capacities.
Away from their tails the sequence is uniformly Cauchy.
The quasicontinuity of each $\widetilde u_j$, with another
summable family of exceptional open sets, makes its limit
energy-quasicontinuous.  This limit agrees almost everywhere
with the $L^1$ limit.

For fixed $k$, the vectors $T_ku_j$ are bounded in $H$ by
(MD5).  Since
\[
 \|T_ku_j-T_ku\|_{L^2(\lambda)}^2
                         \leq2k\|u_j-u\|_{L^1(\lambda)},
\]
their $L^2$ limit is $T_ku$.  Weak compactness in $H$ and
coercivity identify every weak cluster point with this vector.
Lower semicontinuity proves (MD8).
The truncation of the capacity-limit representative is itself
quasicontinuous and agrees $\lambda$-almost everywhere with
this $H$ vector.  Uniqueness of quasicontinuous representatives
therefore identifies it with the canonical representative of
$T_ku$; the capacity estimate consequently concerns the same
representative throughout.
The bound on the capacity of its level sets follows by applying
(MD1) to $|T_ku|/k$ and its canonical representative.
It also shows directly that the possible infinite-value set
has capacity zero.

For $f\in L^\infty(F,\lambda)$, finite-energy symmetry gives
\[
                  \int_Fu_j f\,d\lambda
                      =-\int_F\widetilde{Gf}\,d\mu_j .
\]
The left side converges by $L^1$ convergence and the right
side by total-variation convergence and the bound
$|\widetilde{Gf}|\leq2\|f\|_\infty$.
Thus $u$ satisfies (MD7).  Since bounded measurable functions
separate $L^1$ classes, (MD7) proves uniqueness.
It also proves linearity.  Positivity and order follow by
applying (MD11) to common finite-energy exhaustions of the
positive data and passing to the capacity limit.

We verify the current statement without substituting any
restricted test family.  For original
$\phi\in A_c^{0,0}(V;\mathbf R)$, (AE10) gives
$D(\phi S)=b_\phi\lambda$ with $b_\phi$ bounded and compactly
supported.  Coefficient domination makes $uS$ well defined
on every original complementary form, because $u\in L^1(\lambda)$.
Since $\phi\in H$, the Green identity implies
$G(b_\phi)=-\phi$.  Applying (MD7) with $f=b_\phi$ gives
\[
 \langle D(uS),\phi\rangle
        =\int_Fu\,b_\phi\,d\lambda
        =\int_F\phi\,d\mu .
 \tag{MD12}
\]
This proves (MD9) for all original tests.  Alternatively it
follows by $L^1(\lambda)$ convergence in the same coefficient
pairing from the finite-energy equations.

For two data, choose a common finite-energy exhaustion of
$|\mu|+|\nu|$.  Apply (MD5) to the differences of the
approximations, then pass as above.  This proves all of
(MD10) and independence of the initial exhaustion in both
$L^1$ and the canonical capacity class.
If $\mu$ has finite energy, its Riesz solution already
satisfies (MD7), so uniqueness identifies the two solutions.
\end{proof}

\subsection{The relative BT-capacity consequence}

\begin{proposition}[Independence of the faithful reference]
\label{prop:r11-transposition-reference-independence}
Fix $V$ and the original restriction of $S$ to an ordinary
neighborhood of $\overline V$.  The energy capacity $c_V$,
the class $\mathcal M_{\mathrm{diff}}(F)$, and the canonical
solution and original current in
Theorem~\ref{thm:r11-diffuse-transposition-dirichlet} do not depend
on the choices of compensators, faithful potential $h$, or ambient
affinoid neighborhoods for which the construction is available.
\end{proposition}

\noindent\textbf{Status: \texttt{NEW\_PROVED}} for the larger
measure-data solution class.

\begin{proof}
The original core $A_c^{0,0}(V;\mathbf R)$ and its energy
$\mathcal E(f,g)=\int d'f\wedge d''g\wedge S$ are independent
of all these choices.  Coercivity identifies each part space
with the Hilbert completion of this same core modulo zero
energy, by the same energy norm.  The compact-core formula
for transient capacity therefore gives the same $c_V$ on
compact sets and hence, by outer and Choquet regularity, on
all Borel sets.  In particular all canonical core-limit
quasicontinuous representatives and all polar sets coincide.
The two faithful measures have the same null ideal by (AE7).

Finite-energy Riesz potentials are the same abstract Hilbert
vectors and therefore have the same canonical representatives.
Choose one finite-energy nest for $|\mu|$ relative to this
common energy.  The canonical solution is the limit in the
common capacity of the same finite-energy potential sequence.
Equation (MD11) makes this limit independent of the reference;
its two $L^1$ realizations are representatives of that one
class.  Finally the coefficient multiplication $uS$ and its
distributional derivative use only the original $S$ and
the common almost-everywhere class, proving the current assertion.
\end{proof}

\begin{proposition}[Local BT-capacity control and diffuse BT data]
\label{prop:r11-bt-capacity-diffuse-dirichlet}
For each compact $K_0\Subset V$ there is a finite constant
$C_{K_0}$ depending only on the fixed background, domains and
coefficient controls such that
\[
 \operatorname{Cap}_S^*
   \bigl(K_0\cap\{|\widetilde u_\mu-\widetilde u_\nu|>k\};\Omega\bigr)
       \leq C_{K_0}
           \left(\frac{3\|\mu-\nu\|_{\mathrm{TV}}}{k}\right)^{1/2}.
 \tag{MD13}
\]
The representatives may be set to zero off $\operatorname{supp}S$.
They are locally quasicontinuous for relative BT outer capacity.

In particular every finite signed Radon measure on $V$ whose
variation vanishes on all locally relative BT outer-null sets
is an allowed datum for Theorem~\ref{thm:r11-diffuse-transposition-dirichlet}.
This includes any finite datum with a capacity domination
condition whose right side vanishes at capacity zero; no Dini
growth condition is required for this $L^1$ transposition class.
\end{proposition}

\noindent\textbf{Status: \texttt{DIRECT\_COROLLARY}} of the
new theorem and the frozen one-way energy-to-BT-capacity bridge,
now valid on the analytic affinoid domains of
Theorem~\ref{thm:r11-general-affinoid-faithful-form}.

\begin{proof}
The R4 bridge uses only the faithful coefficient bounds,
local smooth regularization, fixed-cutoff energy estimates,
and the regular compact-core capacity formula.  All these
inputs hold on the present analytic domains by (AE5)--(AE13)
and the frozen local bounded BT calculus; its proof is
unchanged.  Hence
\[
 \operatorname{Cap}_S^*(A;\Omega)
       \leq C_{K_0}\operatorname{cap}_{\mathcal E,1}
                  (A\cap\operatorname{supp}S)^{1/2}
                         \quad(A\subset K_0).
 \tag{MD14}
\]
The full-space energy capacity is no larger than the part
capacity, and (MD2) bounds the latter by $3c_V$.
Equation (MD10) gives (MD13).

Energy quasicontinuity passes through (MD14) on every compact
stage.  The open complement of $\operatorname{supp}S$ has
zero relative capacity, since every relative BT measure is
supported on $\operatorname{supp}S$; it may be included in
each exceptional open set.  This justifies the zero extension
without claiming continuity across its boundary.

An energy-polar set is locally BT outer-null by (MD14).
Thus a datum satisfying the displayed null-set hypothesis
vanishes on it.  The same hypothesis makes it vanish off
$\operatorname{supp}S$.  It belongs to
$\mathcal M_{\mathrm{diff}}(F)$, proving the final assertion.
This argument uses only the null ideals, not equality of the
two capacities or outer regularity of raw BT capacity.
\end{proof}

\subsection{A strict extension beyond finite-energy data}

\begin{proposition}[An original finite diffuse datum of infinite energy norm]
\label{prop:r11-infinite-energy-diffuse-example}
For every $n\geq3$ there are a smooth original positive closed
background $S$ on an ordinary torus box, a finite positive Radon
measure $\mu$ carried by a compact subset, and a part domain $V$,
such that $\mu$ is energy diffuse but
\[
                              \|\mu\|_{-1,V}=\infty.
 \tag{MD15}
\]
Theorem~\ref{thm:r11-diffuse-transposition-dirichlet} supplies
its canonical $L^1$ relative solution; that solution is not in $H$.
\end{proposition}

\noindent\textbf{Status: \texttt{NEW\_PROVED}} as a full-dimensional
strictness example.  It is not used to replace the arbitrary-$S$
theorem by a toric theorem.

\begin{proof}
Use original torus coordinates $t_i=\log|T_i|$,
$\beta=\sum_i d't_i\wedge d''t_i$, and
$S=\beta^{n-1}[U]/(n-1)!$ on a bounded ordinary torus box.
Take a smaller Euclidean ball centered at the origin in its
tropical image and a nonnegative compact smooth cutoff $\chi$
which equals one on a still smaller ball.  Put
\[
 \alpha=(n+2)/2,\qquad
                 \mu=\sigma_*(\chi(t)|t|^{-\alpha}\,dt),
 \tag{MD16}
\]
where $\sigma$ is the original joint Gauss section.
Since $\alpha<n$, this is finite and compactly supported.
The smooth coefficient measure $\beta\wedge S=n\,\sigma_*dt$
is locally dominated by the faithful reference measure.
Consequently $\mu\ll\lambda$, so it charges no energy-polar set.

Let $a=(n-2)/2$, and choose an original compact smooth scalar
$\phi_\varepsilon(t)$ which equals
$(|t|^2+\varepsilon^2)^{-a/2}$ on the smaller ball and has
support in the fixed cutoff ball.  Direct radial integration
with the original calibration gives, for constants independent
of $\varepsilon$,
\[
 \int\phi_\varepsilon\,d\mu\geq c\log(1/\varepsilon)-C,\qquad
             \mathcal E(\phi_\varepsilon,\phi_\varepsilon)
                      \leq C\bigl(1+\log(1/\varepsilon)\bigr).
 \tag{MD17}
\]
Indeed the first integrand has radial exponent
$n-1-a-\alpha=-1$ on $\varepsilon\leq|t|\leq r$,
and the squared-gradient energy has exponent
$n-1-2a-2=-1$ there.  The inner ball and the fixed cutoff
annulus contribute bounded amounts to the energy.
The ratio in the negative-energy core norm tends to infinity,
proving (MD15).  If the resulting solution belonged to $H$,
(AE15) and Cauchy--Schwarz would make that norm finite,
a contradiction.
\end{proof}

\paragraph{The larger class and the remaining nonlinear problem.}
Positive data give positive relative curvature $D(u_\mu S)$,
but $u_\mu$ is not asserted ambient classical psh.  The R10
compact-source obstruction therefore does not conflict with
this theorem.  For an arbitrary original background $S_0$
of degree $(p,p)$ and fixed bounded classical slots
$v_2,\ldots,v_{n-p}$, the entire construction applies to
$S=T_{S_0,n-p-1}(v_2,\ldots,v_{n-p})$.  It solves one free
relative slot on every original test.  It does not construct
a diagonal product of several moving $L^1$ solutions, prove a
general nonlinear Monge--Ampere existence theorem, or settle
OPEN-GQ/OPEN-CTO.
